% (User input window)
% 2026-02-14 10:03:12（Asia/Singapore）

\paragraph{特征标筛—热力学指数率：分拆共振、钩型通道与 Perron 纯性阻碍}
本段在 Schur 断层成像公式（定理 \ref{thm:pom-schur-cycle-index-tomography}）的基础上，对 \(m\to\infty\) 的指数率作一个纯代数的压缩：
将通道迹 \(T_{q,\lambda}(m)=\Tr(B_{q,\lambda}^m)\) 写成按分拆（循环类型）分组的有限线性组合，并用各矩阶压力 \(P(r)\) 的有限组合优化给出通道谱半径的闭式。
在该口径下，钩型通道的“Perron 纯性”（主指数项不允许出现在非平凡通道）要求出现分拆共振并满足一个可机器核验的主项消去约束。

\paragraph{设定：矩阶压力与分拆单项式}
固定 \(q\ge 2\)，并假设对每个 \(1\le r\le q\) 存在原始非负矩阵 \(A_r\) 与向量 \(u_r,v_r\ge 0\) 使得
\[
S_r(m)=u_r^{\top}A_r^{\,m}v_r,
\qquad
\rho_r:=\rho(A_r),
\qquad
P(r):=\log\rho_r.
\]
由 Perron--Frobenius 理论存在常数 \(\kappa_r>0\) 与 \(0<\theta_r<\rho_r\) 使
\begin{equation}\label{eq:pom-sr-main-term-kappa}
S_r(m)=\kappa_r\,e^{mP(r)}+O(\theta_r^{\,m})\qquad(m\to\infty).
\end{equation}
对任意分拆 \(\nu\vdash q\)，记其循环计数为 \(c_r(\nu)\)（即循环类型为 \(1^{c_1}2^{c_2}\cdots\)），并定义分拆单项式
\[
\mathsf{M}_\nu(m):=\prod_{r\ge 1} S_r(m)^{\,c_r(\nu)}\qquad(\text{实际仅需 }1\le r\le q).
\]
并记
\[
z_\nu:=\prod_{r\ge 1} r^{c_r(\nu)}\,c_r(\nu)!\,,
\]
从而共轭类大小为 \(|\mathrm{Cl}(\nu)|=q!/z_\nu\)。

\begin{proposition}[按共轭类压缩的 Schur 断层成像]\label{prop:pom-schur-tomography-by-partitions}
沿用定理 \ref{thm:pom-schur-cycle-index-tomography} 的记号，则对任意 \(\lambda\vdash q\) 与任意 \(m\ge 1\) 有
\[
\boxed{\;
T_{q,\lambda}(m)
=
f^\lambda\sum_{\nu\vdash q}\frac{\chi^\lambda(\nu)}{z_\nu}\,\mathsf{M}_\nu(m)\; }.
\]
\end{proposition}
\begin{proof}
由定理 \ref{thm:pom-schur-cycle-index-tomography}，
\(
T_{q,\lambda}(m)
=\frac{f^\lambda}{q!}\sum_{\sigma\in S_q}\chi^\lambda(\sigma)\,\mathcal{C}_\sigma(m)
\)。
按共轭类分组，并用 \(\mathcal{C}_\sigma(m)=\mathsf{M}_{\nu(\sigma)}(m)\) 与 \(|\mathrm{Cl}(\nu)|=q!/z_\nu\) 得
\[
\sum_{\sigma\in S_q}\chi^\lambda(\sigma)\,\mathcal{C}_\sigma(m)
=\sum_{\nu\vdash q}\frac{q!}{z_\nu}\,\chi^\lambda(\nu)\,\mathsf{M}_\nu(m),
\]
代回即得结论。证毕。
\end{proof}

\paragraph{分拆共振间隙与钩型通道的必要条件}
定义分拆共振间隙（以 \((q)\) 为基准）：
\[
\Delta_{\mathrm{part}}(q)
:=
P(q)-\max_{\substack{\nu\vdash q\\ \nu\neq (q)}}\ \sum_{r\ge 1}c_r(\nu)\,P(r).
\]

\begin{theorem}[钩型通道 \texorpdfstring{$\Rightarrow$}{=>} 共振间隙闭合]\label{thm:pom-hook-channel-implies-partition-gap-closed}
进一步假设 \(A_q\) 原始，且其 Perron 根 \(\rho_q\) 简单并仅出现在平凡通道 \(\lambda=(q)\)（即对所有 \(\lambda\neq(q)\) 有 \(\rho(B_{q,\lambda})<\rho_q\)）。
若存在钩型分拆 \(\lambda=(q-k,1^k)\)（\(1\le k\le q-1\)）使对应通道非零（\(B_{q,\lambda}\neq 0\)），则必有
\[
\boxed{\ \Delta_{\mathrm{part}}(q)=0\ }.
\]
等价地，存在某个非平凡分拆 \(\nu\neq(q)\) 使
\[
P(q)=\sum_{r\ge 1}c_r(\nu)\,P(r).
\]
\end{theorem}
\begin{proof}
由命题 \ref{prop:pom-schur-tomography-by-partitions}，
\[
T_{q,\lambda}(m)
=
f^\lambda\left(\frac{\chi^\lambda((q))}{q}\,S_q(m)+\sum_{\nu\vdash q,\ \nu\neq(q)}\frac{\chi^\lambda(\nu)}{z_\nu}\,\mathsf{M}_\nu(m)\right).
\]
对钩型 \(\lambda=(q-k,1^k)\)，由 Murnaghan--Nakayama 规则知 \(\chi^\lambda((q))=(-1)^k\neq 0\)，故 \(S_q(m)\) 项必出现。
若 \(\Delta_{\mathrm{part}}(q)>0\)，则对一切 \(\nu\neq(q)\) 有
\(
\sum_r c_r(\nu)P(r)\le P(q)-\Delta_{\mathrm{part}}(q)
\)。
由 \eqref{eq:pom-sr-main-term-kappa} 得
\[
\mathsf{M}_\nu(m)
=
\Bigl(\prod_r \kappa_r^{c_r(\nu)}\Bigr)\exp\!\Bigl(m\sum_r c_r(\nu)P(r)\Bigr)\cdot(1+o(1)),
\]
从而所有 \(\nu\neq(q)\) 的项在指数尺度上严格低于 \(S_q(m)\sim \kappa_q e^{mP(q)}\)。
因此
\[
T_{q,\lambda}(m)
=
f^\lambda\frac{\chi^\lambda((q))}{q}\,\kappa_q e^{mP(q)}\cdot(1+o(1)),
\]
推出
\(
\limsup_{m\to\infty}\frac{1}{m}\log|T_{q,\lambda}(m)|=P(q)
\)。
另一方面，对任意有限矩阵 \(B\) 有
\(
\limsup_{m\to\infty}\frac{1}{m}\log|\Tr(B^m)|=\log\rho(B)
\)
（见定理 \ref{thm:pom-schur-channel-radius-character-sieve} 的证明中的引理），故得 \(\rho(B_{q,\lambda})=\rho_q\)，与“Perron 根仅出现在平凡通道”的假设矛盾。
故必有 \(\Delta_{\mathrm{part}}(q)=0\)。证毕。
\end{proof}

\begin{theorem}[严格分拆间隙 \texorpdfstring{$\Rightarrow$}{=>} 钩型通道湮灭]\label{thm:pom-partition-gap-positive-kills-hooks}
在定理 \ref{thm:pom-hook-channel-implies-partition-gap-closed} 的假设下，若 \(\Delta_{\mathrm{part}}(q)>0\)，则对所有非平凡钩型 \(\lambda=(q-k,1^k)\)（\(1\le k\le q-1\)）均有
\[
\boxed{\ B_{q,\lambda}=0\ },
\]
等价地 \(T_{q,\lambda}(m)=0\) 对一切 \(m\ge 1\) 成立。
\end{theorem}
\begin{proof}
取任意钩型 \(\lambda\neq(q)\)。若 \(B_{q,\lambda}\neq 0\)，由定理 \ref{thm:pom-hook-channel-implies-partition-gap-closed} 立得 \(\Delta_{\mathrm{part}}(q)=0\)，与假设矛盾。
故 \(B_{q,\lambda}=0\)，并由 \(T_{q,\lambda}(m)=\Tr(B_{q,\lambda}^m)\) 得等价表述。证毕。
\end{proof}

\paragraph{共振相的主项消去与可判定性}
在 \eqref{eq:pom-sr-main-term-kappa} 的主项制度下，定义共振分拆族
\[
\mathcal{R}_q:=\left\{\nu\vdash q:\ \sum_{r\ge 1}c_r(\nu)\,P(r)=P(q)\right\}.
\]

\begin{theorem}[钩型通道主项消去的代数约束]\label{thm:pom-hook-channel-leading-term-cancellation}
在定理 \ref{thm:pom-hook-channel-implies-partition-gap-closed} 的设定下，若存在非平凡钩型 \(\lambda=(q-k,1^k)\)（\(k\ge 1\)）满足 \(B_{q,\lambda}\neq 0\)，
则必有严格恒等式
\[
\boxed{\;
\frac{\chi^\lambda((q))}{q}\,\kappa_q
+\sum_{\nu\in\mathcal{R}_q\setminus\{(q)\}}
\frac{\chi^\lambda(\nu)}{z_\nu}\,
\prod_{r\ge 1}\kappa_r^{\,c_r(\nu)}
=0\; }.
\]
在给定 \(\{P(r)\}_{1\le r\le q}\) 与 \(\{\kappa_r\}_{1\le r\le q}\) 的指数主阶层面，该恒等式等价于“Perron 纯性”要求：非平凡钩型通道的指数率严格小于 \(P(q)\)。
\end{theorem}
\begin{proof}
由 \eqref{eq:pom-sr-main-term-kappa} 与命题 \ref{prop:pom-schur-tomography-by-partitions}，有渐近展开
\[
\begin{aligned}
T_{q,\lambda}(m)
&=
f^\lambda\sum_{\nu\vdash q}\frac{\chi^\lambda(\nu)}{z_\nu}
\Bigl(\prod_r \kappa_r^{c_r(\nu)}\Bigr)
\exp\!\Bigl(m\sum_r c_r(\nu)P(r)\Bigr)\cdot(1+o(1)).
\end{aligned}
\]
将指数率按大小分层；当 \(\Delta_{\mathrm{part}}(q)=0\) 时，指数率最大的层恰为 \(\mathcal{R}_q\)，且其公共指数为 \(P(q)\)。
因此 \(T_{q,\lambda}(m)\) 的 \(e^{mP(q)}\) 系数（乘去公共因子 \(f^\lambda\)）正是所示和式。
若该系数非零，则 \(\limsup \frac1m\log|T_{q,\lambda}(m)|=P(q)\)，从而 \(\rho(B_{q,\lambda})=\rho_q\)，违背 Perron 纯性假设；
反之，若该系数为零，则最大指数层被精确消去，剩余项指数率严格小于 \(P(q)\)，从而满足 \(\rho(B_{q,\lambda})<\rho_q\) 的指数主阶必要条件。
证毕。
\end{proof}

\begin{theorem}[分拆共振的比值判别（可核验）]\label{thm:pom-partition-resonance-ratio-test}
对任意 \(q\ge 2\) 与任意分拆 \(\nu\vdash q\)，令 \(\mathsf{M}_\nu(m)=\prod_r S_r(m)^{c_r(\nu)}\)。
则在上述主项制度下有
\[
\boxed{\;
\limsup_{m\to\infty}\frac{1}{m}\log\frac{S_q(m)}{\mathsf{M}_\nu(m)}
=P(q)-\sum_{r\ge 1}c_r(\nu)\,P(r)\; }.
\]
特别地，
\[
P(q)=\sum_{r}c_r(\nu)P(r)
\quad\Longleftrightarrow\quad
\lim_{m\to\infty}\frac{S_q(m)}{\mathsf{M}_\nu(m)}
=\frac{\kappa_q}{\prod_{r}\kappa_r^{c_r(\nu)}}\in(0,\infty).
\]
\end{theorem}
\begin{proof}
由 \eqref{eq:pom-sr-main-term-kappa} 得
\(
\frac{1}{m}\log S_r(m)\to P(r)
\)
并且
\(
\frac{S_r(m)}{\kappa_r e^{mP(r)}}\to 1
\)。
因此
\[
\frac{1}{m}\log \mathsf{M}_\nu(m)
=\sum_r c_r(\nu)\cdot \frac{1}{m}\log S_r(m)\ \longrightarrow\ \sum_r c_r(\nu)P(r),
\]
从而
\(
\limsup\frac1m\log\frac{S_q(m)}{\mathsf{M}_\nu(m)}
=P(q)-\sum_r c_r(\nu)P(r)
\)。
若再用主项比例极限 \(S_r(m)\sim \kappa_r e^{mP(r)}\)，则比值极限的常数亦立即得到。证毕。
\end{proof}

\paragraph{Schur 通道谱半径的“特征标筛—热力学”公式}
对任意 \(\lambda\vdash q\)，定义特征标支持集
\[
\mathrm{Supp}(\lambda):=\{\nu\vdash q:\ \chi^\lambda(\nu)\neq 0\},
\]
以及筛后压力
\[
\mathcal{P}_\lambda:=\max_{\nu\in\mathrm{Supp}(\lambda)}\ \sum_{r\ge 1}c_r(\nu)\,P(r).
\]
并令极大值集合
\[
\mathcal{M}_\lambda:=\left\{\nu\in\mathrm{Supp}(\lambda):\ \sum_{r}c_r(\nu)P(r)=\mathcal{P}_\lambda\right\}.
\]

\begin{theorem}[通道谱半径的特征标筛公式（非退化主项支配）]\label{thm:pom-schur-channel-radius-character-sieve}
在上述设定下，假设极大值集合上的主系数加权和非零：
\[
\sum_{\nu\in\mathcal{M}_\lambda}\frac{\chi^\lambda(\nu)}{z_\nu}\prod_{r\ge 1}\kappa_r^{c_r(\nu)}\neq 0.
\]
则通道谱半径满足严格等式
\[
\boxed{\ \rho(B_{q,\lambda})=\exp(\mathcal{P}_\lambda)\ }.
\]
从而 near-RH 型平方根门槛在该通道上等价于纯标量不等式
\[
\boxed{\ 2\mathcal{P}_\lambda\le P(q)\ }.
\]
\end{theorem}
\begin{proof}
由命题 \ref{prop:pom-schur-tomography-by-partitions} 与主项展开 \eqref{eq:pom-sr-main-term-kappa}，
\[
T_{q,\lambda}(m)
=
f^\lambda\sum_{\nu\vdash q}\frac{\chi^\lambda(\nu)}{z_\nu}
\Bigl(\prod_r \kappa_r^{c_r(\nu)}\Bigr)
\exp\!\Bigl(m\sum_r c_r(\nu)P(r)\Bigr)\cdot(1+o(1)).
\]
在“非退化”假设下，指数率最大的 \(\mathcal{P}_\lambda\) 层不发生完全消去，因此
\(
\limsup_{m\to\infty}\frac{1}{m}\log|T_{q,\lambda}(m)|=\mathcal{P}_\lambda
\)。
另一方面，对任意有限维矩阵 \(B\) 记其特征值为 \(\{\mu_j\}\)，则
\(
\Tr(B^m)=\sum_j \mu_j^m
\)。
令 \(\rho=\rho(B)=\max_j|\mu_j|\)，并记 \(\eta_j=\mu_j/\rho\)（使得 \(|\eta_j|\le 1\) 且至少一个满足 \(|\eta_j|=1\)），则
\[
\Tr(B^m)=\rho^m\Bigl(\sum_{|\eta_j|=1}\eta_j^m+O((\rho'/\rho)^m)\Bigr),
\]
其中 \(0\le \rho'<\rho\) 为次级谱半径。与定理 \ref{thm:finite-part-cyclic-lift-limsup-spectral-gap} 中同型的 Ces\`aro 平方模论证可得
\(\limsup_{m\to\infty}\bigl|\sum_{|\eta_j|=1}\eta_j^m\bigr|>0\)，
从而
\[
\limsup_{m\to\infty}\frac{1}{m}\log|\Tr(B^m)|=\log\rho(B).
\]
应用于 \(B=B_{q,\lambda}\) 并与上一段的指数率识别合并，即得 \(\log\rho(B_{q,\lambda})=\mathcal{P}_\lambda\)。
最后，\(2\mathcal{P}_\lambda\le P(q)\) 等价于 \(\rho(B_{q,\lambda})^2\le \rho_q\)（其中 \(\rho_q=e^{P(q)}\)），即平方根门槛。证毕。
\end{proof}

